\begin{table}[t]
\centering
\caption{Block space of the anchor mechanism per force close, observed (two P2WSH anchors and their spends) and under zero-fee commitments (one pay-to-anchor output spent by a child).}
\label{tab:zero_fee}
\footnotesize
\begin{tabular}{lrrrrr}
\toprule
Year & Anchor closes & Bumped \% & Anchors vB & P2A vB & Anchor BTC \\
\midrule
2021 & 3{,}839 & 8.4 & 194.1 & 155.9 & 0.022 \\
2022 & 34{,}297 & 12.2 & 200.4 & 151.8 & 0.207 \\
2023 & 39{,}280 & 36.6 & 180.5 & 124.6 & 0.231 \\
2024 & 41{,}035 & 12.6 & 163.1 & 151.3 & 0.226 \\
2025 & 59{,}927 & 2.9 & 166.0 & 162.0 & 0.339 \\
2026\textsuperscript{a} & 45{,}209 & 4.1 & 162.0 & 160.7 & 0.253 \\
\bottomrule
\end{tabular}
\par\smallskip\parbox{\linewidth}{\scriptsize \textsuperscript{a}Closes to 4 October 2026. Zero-fee closes that were not already bumped need a child of their own (wallet input and change output, 111.0 vB). Anchor BTC: value locked in 330-sat anchors, which pay-to-anchor outputs (0 sat) do not create.}
\end{table}
